\documentclass[10pt,a4paper]{amsart}
\usepackage[margin=19mm]{geometry}
\usepackage{amsmath,amssymb,mathtools}
\usepackage[scr=rsfs,cal=euler]{mathalfa}
\usepackage{xfrac}
\usepackage[unicode,hidelinks,bookmarks=false]{hyperref}
\newcommand{\ie}{i.e.\ }
\newcommand{\cf}{cf.\ }
\newcommand{\resp}{respectively\ }
\newcommand{\Subsection}{\subsection}
\providecommand{\nobreakdash}{\nobreak\hbox{-}}
\setlength{\emergencystretch}{2em}
\multlinegap0pt
\usepackage{tikz}
\usepackage{amssymb}
\newtheorem{thm}{Theorem}
\newtheorem{prop}{Proposition}
\newtheorem{lem}{Lemma}
\DeclareMathOperator{\I}{\mathtt{I}}
\DeclareMathOperator{\II}{\mathtt{II}}
\DeclareMathOperator{\III}{\mathtt{III}}
\newcommand{\Q}{\mathbf{Q}}
\newcommand{\T}{\operatorname{\mathsf{T}}}
\newcommand{\Z}{\mathbf{Z}}
\newcommand{\tors}{\mathsf{tors}}
\DeclareMathOperator{\typ}{{typ}}
\title{Kodaira-Neron statistics for rational elliptic curves with $j$-invariant 0 and 1728}
\author{John Cullinan, Sebastian Sargenti}
\date{Source: arXiv:2605.14226v1 (CC BY 4.0)}
\begin{document}
\maketitle

\noindent\emph{Source and licence.} Kodaira-Neron statistics for rational elliptic curves with $j$-invariant 0 and 1728. Version arXiv:2605.14226v1 (CC BY 4.0), distributed by arXiv under the Creative Commons Attribution 4.0 International licence, http://creativecommons.org/licenses/by/4.0/, as recorded in the arXiv licence metadata for this exact version and shown on the abstract page https://arxiv.org/abs/2605.14226v1. Full licence notice: \texttt{licenses/arxiv-2605.14226-CC-BY-4.0.txt}.\par\smallskip

\noindent\textit{Editorial scope.} Counted unit: the article's thm j1728\_thm (source0.tex lines 116--142). The article's complete original proof of it is delivered with it (source0.tex lines 361--411, one \texttt{proof} environment of 51 lines, which the article places away from the statement). No mathematics is written, repaired or re-derived: the statement and the proof are the article's own lines, in the article's own notation, and no retained line is substituted or re-flowed.  Retained uncounted as the source-local context the proof needs: lines 218--220; lines 287--301; lines 340--349; lines 639--656; lines 658--670; lines 672--684; lines 686--702.  Nothing was omitted from any retained span, so the package carries no omission record.  No retained line contains a citation, so the wrapper carries no bibliography and no bibliography entry.  No retained line contains a figure environment, an image-inclusion command or drawing code, so no image is delivered and the package declares no asset dependency.  Editorial formatting only, in addition: an AMS \texttt{amsart} preamble that restates the article's own declarations and macros used by the retained lines, namely the environments \texttt{thm}, \texttt{prop}, \texttt{lem} on separate counters, together with the standard packages those lines need.  The retained environments print in this document as Theorem 1, Proposition 1, Lemma 1 in order of appearance, whereas the article prints the same items with the numbering of its own full document.  The complete native source of the article as distributed in the arXiv e-print for this exact version is bundled unchanged as \texttt{sources/200567/source0.tex}.
\begin{thm} \label{1728_min_thm}
Let $E$ be an elliptic curve over $\Q$ with $j$-invariant 1728 written in short Weierstrass form $y^2 = x^3 + Ax$ with $A \in \Z$. Then $E$ is minimal if and only if  $A$ is fourth-power-free.  Moreover, $\typ_2(E)$ is given in Tables \ref{j1728_finite_tables}, \ref{j1728_(2,2)tors},  \ref{j1728-Z2tors-size4}, and \ref{j1728-generic}. 
\end{thm}

\begin{prop} \label{1728_minimal_models}
Let $E/\Q$ be a minimal elliptic curve with $j$-invariant 1728.  Then $E$ is isomorphic to \href{https://beta.lmfdb.org/EllipticCurve/Q/32/a/3}{32.a3}, \href{https://beta.lmfdb.org/EllipticCurve/Q/32/a/4}{32.a4},  \href{https://beta.lmfdb.org/EllipticCurve/Q/64/a/3}{64.a3}, \href{https://beta.lmfdb.org/EllipticCurve/Q/64/a/4}{64.a4}, or is a nonzero integral specialization of one of the following curves:
\begin{center}
\begin{tabular}{|c|c|l|l|}
\hline
Graph & $E(\Q)_\tors$ & Model &  Conditions\\
\hline
$T_4^3$ & $\Z/2\Z \times \Z/2\Z$ & $y^2 = x^3-t^2x$ & $t$ squarefree\\
& $\Z/2\Z$ &$y^2 = x^3+t^2x$ & $t$ squarefree \\
\hline
$L_2(2)$ & $\Z/2\Z$ & $y^2 = x^3 +tx$ & $t \ne \pm s^2$, $t$ 4th-power-free \\
\hline
\end{tabular}
\end{center}
\end{prop}

\begin{lem} \label{squarefree}
The proportion of $k$th-power-free positive integers $n$ such that:
\begin{enumerate}
\item $k=2$ and $n$ is odd is $(2/3) \cdot (1/\zeta(2))$; \label{odd_sqf}
\item $k=2$ and $n$ is coprime to 3, is $(3/4) \cdot (1/\zeta(2))$;
\item $k=3$ and $n$ is odd is $(4/7) \cdot (1/\zeta(3))$;
\item $k=3$ and $v_2(n)=1$ is $(2/7) \cdot (1/\zeta(3))$;
\item $k=4$ and $n$ is odd is $(8/15) \cdot (1/\zeta(4))$;
\end{enumerate} 
\end{lem}

\begin{center}
\begin{table}[h] 
\begin{tabular}{|cl|}
\hline
$E$ & $\typ_2(E)$ \\
\hline
\href{https://beta.lmfdb.org/EllipticCurve/Q/32/a/3}{32.a3} & $\III$  \\
\href{https://beta.lmfdb.org/EllipticCurve/Q/34/a/4}{32.a4} & $\I_3^*$ \\
\href{https://beta.lmfdb.org/EllipticCurve/Q/64/a/3}{64.a3} & $\I_2^*$ \\
\href{https://beta.lmfdb.org/EllipticCurve/Q/64/a/4}{64.a4} & $\II$ \\
\hline
\end{tabular}
\caption{$\typ_2(E)$ for \href{https://beta.lmfdb.org/EllipticCurve/Q/32/a/3}{32.a3},
\href{https://beta.lmfdb.org/EllipticCurve/Q/32/a/4}{32.a4}, 
\href{https://beta.lmfdb.org/EllipticCurve/Q/64/a/3}{64.a3}, 
\href{https://beta.lmfdb.org/EllipticCurve/Q/64/a/4}{64.a4}} \label{j1728_finite_tables}
\end{table} 
\end{center}

\begin{center}
 \begin{table}[h]
 \begin{tabular}{|l|l|}
 \hline
$\typ_2(E)$ & Conditions \\
 \hline
 $\I_2^*$ & $v_2(a) = 2$\\
 $\III$ & $v_2(a) = 0$\\ 
  \hline 
\end{tabular}
 \caption{$j$-invariant 1728, $E(\Q)_\tors \simeq \Z/2\Z \times \Z/2\Z$} \label{j1728_(2,2)tors}
 \end{table}
 \end{center}

\begin{center}
\begin{table}[h]
 \begin{tabular}{|l|l|}
 \hline
$\typ_2(E)$ & Conditions \\
 \hline
 $\II$ & $t$ even \\
 $\I_3^*$ & $t$ odd \\
\hline
\end{tabular}
 \caption{$j$-invariant 1728, $E(\Q)_\tors \simeq \Z/2\Z$, Isogeny Class size 4} \label{j1728-Z2tors-size4}
 \end{table}
  \end{center}

\begin{center}
 \begin{table}[h]
 \begin{tabular}{|l|l|}
 \hline
$\typ_2(E)$ & Conditions \\
 \hline
 $\I_2^*$ & $v_2(a)=2$ and $a/4 \equiv 3 \pmod{4}$ \\
$\I_3^*$ & $v_2(a) = 2$ and $a/4 \equiv 1 \pmod{4}$\\
$\II$ & $v_2(a)=0$ and $a \equiv 1 \pmod{4}$\\
$\III$ & $v_2(a)=1$, or \\
 &$v_2(a)=0$ and $a \equiv 3 \pmod{4}$\\
 $\III^*$ & $v_2(a)=3$ \\
\hline
\end{tabular}
 \caption{$j$-invariant 1728, $E(\Q)_\tors \simeq \Z/2\Z$, Isogeny Class size 2} \label{j1728-generic}
 \end{table}
 \end{center}

\begin{thm} \label{j1728_thm}
Let $E$ be an elliptic curve over $\Q$ with $j$-invariant 1728.  Let $\T = \typ_2(E)$.  
\begin{enumerate}
\item If $E$ belongs to an isogeny-torsion graph isomorphic to $T_4^3$, then \\
$\T \in \lbrace \II, \III, \I_2^*, \I_3^* \rbrace$ and 
\[
N_{T_4^3,\T}^{1728}(X) = \frac{c_{\T}}{\zeta(2)\sqrt[3]{2}}X^{1/6} + O(X^{1/12}),
\]
where
\[
c_{\T} = \begin{cases} 2/3 & \text{ if } \T =   \II \text{ or }\III, \\
1/3 & \text{ if } \T = \I_2^* \text{ or } \I_3^*. \end{cases}
\]
\item If $E$ belongs to an isogeny-torsion graph isomorphic to $L_2(2)$, then \\
$\T \in \lbrace \II, \III, \III^*, \I_2^*, \I_3^* \rbrace$ and 
and 
\begin{align*}
N_{L_2(2),\T}^{1728}(X) =  \frac{c_{\T}}{\zeta(4)\sqrt[3]{4}}X^{1/3} + O(X^{1/6}),
\end{align*}
where
\begin{align} \label{c_T_first_time}
c_{\T} = \begin{cases} 8/15 & \text{ if } \T = \III, \\
4/15 & \text{ if } \T = \II, \\
1/15 & \text{ if } \T = \I_2^*,  \I_3^*, \text{ or } \III^*. \end{cases}
\end{align}
\end{enumerate}
\end{thm} 

\begin{proof}[Proof of Theorem \ref{j1728_thm}]
Set $G = T_4^3$ and let $X >0$.  We first count minimal elliptic curves with torsion subgroup $\Z/2\Z \times \Z/2\Z$ and height $\leq X$.  By  Proposition \ref{1728_minimal_models}, these are given by 
\[
y^2 = x^3 -t^2 x,
\]
with $4t^6 <X$ and $t$ squarefree; the number of such integers is 
\[
\frac{1}{\zeta(2)} \cdot \left(\frac{X}{4}\right)^{1/6} + O(X^{1/12}).
\]
Of these elliptic curves, $\typ_2(E) = \III$ if $t$ is odd and $\I_2^*$ if $t$ is even, by Theorem \ref{1728_min_thm} and Table \ref{j1728_(2,2)tors}.  By Lemma \ref{squarefree} the proportion of odd squarefree integers is $2/3\zeta(2)$, while the proportion of even squarefree integers is $1/3\zeta(2)$.  Therefore, 
\[
N_{T_4^3,\T}^{1728}(X) =\frac{c_{\T}}{\zeta(2)\sqrt[3]{2}}X^{1/6} + O(X^{1/12}),
\]
where
\[
c_{\T} = \begin{cases} 2/3 & \text{ if } \T =   \III, \text{ and} \\
1/3 & \text{ if } \T = \I_2^*. \end{cases}
\]
It now remains to determine the asymptotic for the case of $j=1728$, $G = T_4^3$, and torsion subgroup $\Z/2\Z$.  By Proposition \ref{1728_min_thm}, all such minimal elliptic curves are given by 
\begin{align} \label{min_j1728_Z2_tors}
y^2 = x^3 + t^2x,
\end{align}
with $t$ squarefree; by Table  \ref{j1728-Z2tors-size4} they satisfy $\typ_2(E) = \II$ if $t$ is odd and squarefree and $\typ_2(E) = \I_3^*$ if $t$ is even and squarefree.   The height of an elliptic curve with Weierstrass model given by Equation (\ref{min_j1728_Z2_tors}) is $4t^6$.  Therefore, 
\[
N_{T_4^3,\T}^{1728}(X) =\frac{c_{\T}}{\zeta(2)\sqrt[3]{2}}X^{1/6} + O(X^{1/12}),
\]
where
\[
c_{\T} = \begin{cases} 2/3 & \text{ if } \T =   \II, \text{ and} \\
1/3 & \text{ if } \T = \I_3^*. \end{cases}
\]
This completes the proof of the case $G= T_4^3$.  

Next, let $G = L_2(2)$.  Working modulo $2^4$, it follows directly from the arithmetic criteria in Table \ref{j1728-generic} that Kodaira-N\'eron types $\III, \II, \I_2^*, \I_3^*$, and $\III^*$ occur with proportions $8/15$, $4/15$, $1/15$, $1/15$, and $1/15$, respectively.  The models of $j$-invariant 1728 elliptic curves that are minimal,  belong to the graph $L_2(2)$, and have height $\leq X$  are parameterized by fourth-power-free integers $t$ that are not squares such that $4|t|^3\leq X$.

The number of  fourth--power-free integers $t$ with $t \leq (X/4)^{1/3}$ is  
\[
\frac{1}{\zeta(4)} \cdot \left(X/4\right)^{1/3}  + O(X^{1/12}),
\]
while the number of square integers $t \leq  (X/4)^{1/3}$ is  $\sim \sqrt{(X/4)^{1/3}} = O(X^{1/6})$.  It follows that 
\[
N_{L_2(2),\T}^{1728}(X) = \frac{1}{\zeta(4)} \left(\frac{X}{4}\right)^{1/3} + O(X^{1/6}),
\]
where 
\[
c_{\T} = \begin{cases} 8/15 & \text{ if } \T = \III, \\
4/15 & \text{ if } \T = \II, \\
1/15 & \text{ if } \T = \I_2^*,  \I_3^*, \text{ or } \III^*,\end{cases}
\]
as claimed. 
\end{proof}

\end{document}
